\clearpage

\section{Независимое структурное расширение}

Расширение содержит 288 независимых обучающих наборов и 3168 проверенных калибровочных/валидационных эпизодов. Все наборы удовлетворяют заданному предкризисному условию. Известная и обнаруженная структуры используют одинаковые обученные поведенческие параметры и заданные отображения ресурсов. Ручного исправления discovery и выбора политик по тесту нет. Полное машиночитаемое сравнение, включая каждое фиксированное действие и абляции компонентов, находится в архиве: \texttt{analysis/full\_summary.csv}.

\subsection{Стоимость по всем структурам и механизмам}

Таблицы показывают среднее отношений стоимости каждого обученного варианта к соответствующей стоимости без вмешательства, равной 1. Regret — разность retained-load и планировщика с генератором в этой шкале. Частота тяжёлых событий относится к retained-load. K означает известную, D — обнаруженную структуру. В каждой строке 12 независимых обучений; эвристика и генератор одинаковы в обоих структурных вариантах.

\begin{table}[htbp]\caption{P1: нормированная стоимость и частота событий.}\centering\small\setlength{\tabcolsep}{3pt}

\begin{tabular}{llrrrrrrr}\toprule

Model & Mechanism & Static & Pred. & Ret. & Heur. & Gen. & Regret & Event\\\midrule

K & fixed & 0.817 & 0.800 & 0.807 & 0.801 & 0.798 & 0.009 & 0.082\\

K & service & 0.599 & 0.630 & 0.639 & 0.670 & 0.578 & 0.061 & 0.399\\

K & routing & 0.587 & 0.549 & 0.502 & 0.501 & 0.429 & 0.073 & 0.365\\

K & combined & 0.729 & 0.692 & 0.624 & 0.636 & 0.588 & 0.036 & 0.641\\

D & fixed & 0.760 & 0.760 & 0.760 & 0.801 & 0.798 & -0.038 & 0.035\\

D & service & 0.580 & 0.580 & 0.577 & 0.670 & 0.578 & -0.001 & 0.366\\

D & routing & 0.575 & 0.575 & 0.568 & 0.501 & 0.429 & 0.139 & 0.378\\

D & combined & 0.722 & 0.722 & 0.721 & 0.636 & 0.588 & 0.133 & 0.648\\

\bottomrule\end{tabular}\end{table}

\begin{table}[htbp]\caption{P2: нормированная стоимость и частота событий.}\centering\small\setlength{\tabcolsep}{3pt}

\begin{tabular}{llrrrrrrr}\toprule

Model & Mechanism & Static & Pred. & Ret. & Heur. & Gen. & Regret & Event\\\midrule

K & fixed & 0.955 & 0.955 & 0.988 & 1.077 & 0.953 & 0.035 & 0.102\\

K & service & 0.647 & 0.647 & 0.701 & 0.820 & 0.630 & 0.072 & 0.333\\

K & routing & 0.731 & 0.614 & 0.428 & 0.514 & 0.428 & 0.000 & 0.325\\

K & combined & 0.751 & 0.640 & 0.582 & 0.605 & 0.521 & 0.061 & 0.514\\

D & fixed & 0.918 & 0.918 & 1.072 & 1.077 & 0.953 & 0.119 & 0.095\\

D & service & 0.637 & 0.637 & 0.878 & 0.820 & 0.630 & 0.249 & 0.351\\

D & routing & 0.772 & 0.772 & 0.897 & 0.514 & 0.428 & 0.469 & 0.370\\

D & combined & 0.748 & 0.751 & 0.798 & 0.605 & 0.521 & 0.278 & 0.576\\

\bottomrule\end{tabular}\end{table}

\begin{table}[htbp]\caption{P3: нормированная стоимость и частота событий.}\centering\small\setlength{\tabcolsep}{3pt}

\begin{tabular}{llrrrrrrr}\toprule

Model & Mechanism & Static & Pred. & Ret. & Heur. & Gen. & Regret & Event\\\midrule

K & fixed & 0.930 & 0.930 & 0.958 & 0.920 & 0.923 & 0.036 & 0.313\\

K & service & 0.535 & 0.535 & 0.643 & 0.535 & 0.535 & 0.108 & 0.535\\

K & routing & 0.630 & 0.630 & 0.650 & 0.630 & 0.516 & 0.133 & 0.651\\

K & combined & 0.878 & 0.878 & 0.871 & 0.880 & 0.839 & 0.032 & 0.667\\

D & fixed & 0.930 & 0.930 & 0.958 & 0.920 & 0.923 & 0.036 & 0.313\\

D & service & 0.535 & 0.535 & 0.643 & 0.535 & 0.535 & 0.108 & 0.535\\

D & routing & 0.630 & 0.630 & 0.650 & 0.630 & 0.516 & 0.133 & 0.651\\

D & combined & 0.878 & 0.878 & 0.871 & 0.880 & 0.839 & 0.032 & 0.667\\

\bottomrule\end{tabular}\end{table}

\begin{table}[htbp]\caption{P4: нормированная стоимость и частота событий.}\centering\small\setlength{\tabcolsep}{3pt}

\begin{tabular}{llrrrrrrr}\toprule

Model & Mechanism & Static & Pred. & Ret. & Heur. & Gen. & Regret & Event\\\midrule

K & fixed & 0.720 & 0.720 & 0.805 & 0.888 & 0.720 & 0.085 & 0.255\\

K & service & 0.545 & 0.545 & 0.649 & 0.649 & 0.544 & 0.105 & 0.564\\

K & routing & 0.520 & 0.528 & 0.610 & 0.620 & 0.512 & 0.098 & 0.490\\

K & combined & 0.801 & 0.809 & 0.828 & 0.839 & 0.800 & 0.029 & 0.672\\

D & fixed & 0.720 & 0.720 & 0.805 & 0.888 & 0.720 & 0.085 & 0.255\\

D & service & 0.545 & 0.545 & 0.649 & 0.649 & 0.544 & 0.105 & 0.564\\

D & routing & 0.520 & 0.528 & 0.610 & 0.620 & 0.512 & 0.098 & 0.490\\

D & combined & 0.801 & 0.809 & 0.828 & 0.839 & 0.800 & 0.029 & 0.672\\

\bottomrule\end{tabular}\end{table}

\begin{table}[htbp]\caption{P5: нормированная стоимость и частота событий.}\centering\small\setlength{\tabcolsep}{3pt}

\begin{tabular}{llrrrrrrr}\toprule

Model & Mechanism & Static & Pred. & Ret. & Heur. & Gen. & Regret & Event\\\midrule

K & fixed & 0.799 & 0.799 & 0.950 & 0.985 & 0.800 & 0.150 & 0.306\\

K & service & 0.348 & 0.350 & 0.352 & 0.460 & 0.348 & 0.003 & 0.483\\

K & routing & 1.030 & 0.727 & 0.473 & 0.422 & 0.422 & 0.050 & 0.354\\

K & combined & 1.142 & 1.022 & 1.008 & 0.997 & 0.984 & 0.025 & 0.672\\

D & fixed & 0.799 & 0.799 & 0.950 & 0.985 & 0.800 & 0.150 & 0.306\\

D & service & 0.348 & 0.350 & 0.352 & 0.460 & 0.348 & 0.003 & 0.483\\

D & routing & 1.030 & 0.727 & 0.473 & 0.422 & 0.422 & 0.050 & 0.354\\

D & combined & 1.142 & 1.022 & 1.008 & 0.997 & 0.984 & 0.025 & 0.672\\

\bottomrule\end{tabular}\end{table}

\begin{table}[htbp]\caption{P6: нормированная стоимость и частота событий.}\centering\small\setlength{\tabcolsep}{3pt}

\begin{tabular}{llrrrrrrr}\toprule

Model & Mechanism & Static & Pred. & Ret. & Heur. & Gen. & Regret & Event\\\midrule

K & fixed & 0.753 & 0.753 & 0.799 & 1.051 & 0.753 & 0.045 & 0.217\\

K & service & 0.494 & 0.494 & 0.633 & 0.920 & 0.492 & 0.140 & 0.439\\

K & routing & 0.575 & 0.575 & 0.422 & 0.460 & 0.349 & 0.073 & 0.401\\

K & combined & 0.745 & 0.693 & 0.602 & 0.685 & 0.590 & 0.012 & 0.665\\

D & fixed & 0.740 & 0.740 & 0.740 & 1.051 & 0.753 & -0.014 & 0.174\\

D & service & 0.492 & 0.492 & 0.492 & 0.920 & 0.492 & 0.000 & 0.359\\

D & routing & 0.568 & 0.568 & 0.605 & 0.460 & 0.349 & 0.256 & 0.528\\

D & combined & 0.743 & 0.743 & 0.743 & 0.685 & 0.590 & 0.153 & 0.667\\

\bottomrule\end{tabular}\end{table}

\clearpage

\subsection{Основные сравнения и проверка обнаружения структуры}

Шесть основных сравнений при нагрузочных возвратах используют интервалы с поправкой Бонферрони, контролирующей ошибку семейства на уровне 0.05. Ниже приведены абсолютные разности стоимости retained минус predictive, делённые на 1000, а не проценты. Положительная разность означает преимущество предиктивного отбора.

\begin{center}\small\begin{tabular}{lrrr}\toprule

Process & Difference & Lower & Upper\\\midrule

P1 & -1.081 & -2.317 & 0.155\\

P2 & -4.694 & -10.293 & 0.905\\

P3 & 2.056 & -9.309 & 13.420\\

P4 & 5.092 & 1.931 & 8.252\\

P5 & -14.544 & -31.191 & 2.103\\

P6 & -11.514 & -25.155 & 2.126\\

\bottomrule\end{tabular}\end{center}

\begin{center}\small\begin{tabular}{lrrrrrr}\toprule

Process & Fitness & Precision & Equal & Added & Missing & Same policy\\\midrule

P1 & 1.000 & 0.497 & 0/48 & 48/48 & 0/48 & 25/48\\

P2 & 1.000 & 0.642 & 0/48 & 48/48 & 0/48 & 3/48\\

P3 & 1.000 & 0.996 & 48/48 & 0/48 & 0/48 & 48/48\\

P4 & 1.000 & 1.000 & 48/48 & 0/48 & 0/48 & 48/48\\

P5 & 1.000 & 1.000 & 48/48 & 0/48 & 0/48 & 48/48\\

P6 & 1.000 & 0.487 & 0/48 & 48/48 & 0/48 & 17/48\\

\bottomrule\end{tabular}\end{center}

Проверка языков сравнивает видимые последовательности исполнения с фиксированным ограничением числа состояний; незавершённые проверки сохранены в \texttt{discovery\_summary.json}. Same policy считает полное совпадение действия и порога retained-load между известной и обнаруженной структурой во всех четырёх механизмах. 144 модели с добавленными путями зарегистрированы как нарушения соответствия исходной структуре в \texttt{structural\_conformance.json}; их стоимость без исправления сохраняется для анализа ошибок. Этот признак оценки по известному оригиналу отделён от вычислительного отказа и не является реализованным правилом допуска политики.

\subsection{Чувствительность селектора}

Перебор использует сохранённые планировочные решения и независимые тестовые действия. Это диагностика; порог не выбирается по тестовой стоимости. Таблицы относятся к нагрузочным возвратам и известной структуре. D и R — число сохранённых компонентов из 12. Agree — совпадение типа действия с планировщиком по генератору; regret — абсолютная разность стоимости, делённая на 1000. Остальные механизмы и обнаруженные структуры включены в \texttt{selector\_summary.csv}.

\begingroup\small\setlength{\tabcolsep}{3pt}\begin{longtable}{llrrrrrr}\caption{Перебор селектора: return.}\\\toprule

Process & Threshold & D & R & Agree & Cost ratio & Regret & Event\\\midrule\endfirsthead

\toprule Process & Threshold & D & R & Agree & Cost ratio & Regret & Event\\\midrule\endhead

P1 & 0 & 0 & 6 & 0.417 & 0.527 & 2.450 & 0.378\\

P1 & 0.001 & 0 & 3 & 0.250 & 0.549 & 2.920 & 0.384\\

P1 & 0.0025 & 0 & 3 & 0.250 & 0.549 & 2.920 & 0.384\\

P1 & 0.005 & 0 & 3 & 0.250 & 0.549 & 2.920 & 0.384\\

P1 & 0.01 & 0 & 2 & 0.167 & 0.561 & 3.186 & 0.385\\

P2 & 0 & 0 & 8 & 0.667 & 0.566 & 3.540 & 0.340\\

P2 & 0.001 & 0 & 8 & 0.667 & 0.566 & 3.540 & 0.340\\

P2 & 0.0025 & 0 & 8 & 0.667 & 0.566 & 3.540 & 0.340\\

P2 & 0.005 & 0 & 6 & 0.583 & 0.614 & 4.694 & 0.349\\

P2 & 0.01 & 0 & 4 & 0.500 & 0.646 & 5.616 & 0.358\\

P3 & 0 & 0 & 7 & 0.083 & 0.693 & 19.309 & 0.651\\

P3 & 0.001 & 0 & 4 & 0.083 & 0.672 & 17.170 & 0.653\\

P3 & 0.0025 & 0 & 0 & 0.250 & 0.630 & 12.473 & 0.651\\

P3 & 0.005 & 0 & 0 & 0.250 & 0.630 & 12.473 & 0.651\\

P3 & 0.01 & 0 & 0 & 0.250 & 0.630 & 12.473 & 0.651\\

P4 & 0 & 0 & 7 & 0.500 & 0.577 & 3.970 & 0.467\\

P4 & 0.001 & 0 & 7 & 0.500 & 0.577 & 3.970 & 0.467\\

P4 & 0.0025 & 0 & 4 & 0.667 & 0.559 & 2.814 & 0.460\\

P4 & 0.005 & 0 & 1 & 0.917 & 0.528 & 0.953 & 0.436\\

P4 & 0.01 & 0 & 0 & 1.000 & 0.520 & 0.494 & 0.434\\

P5 & 0 & 0 & 9 & 0.750 & 0.575 & 8.443 & 0.378\\

P5 & 0.001 & 0 & 8 & 0.667 & 0.625 & 11.296 & 0.399\\

P5 & 0.0025 & 0 & 6 & 0.500 & 0.727 & 17.412 & 0.427\\

P5 & 0.005 & 0 & 6 & 0.500 & 0.727 & 17.412 & 0.427\\

P5 & 0.01 & 0 & 3 & 0.250 & 0.879 & 26.373 & 0.474\\

P6 & 0 & 0 & 3 & 0.250 & 0.517 & 12.708 & 0.488\\

P6 & 0.001 & 0 & 1 & 0.083 & 0.556 & 15.705 & 0.521\\

P6 & 0.0025 & 0 & 1 & 0.083 & 0.556 & 15.705 & 0.521\\

P6 & 0.005 & 0 & 0 & 0.000 & 0.575 & 17.128 & 0.528\\

P6 & 0.01 & 0 & 0 & 0.000 & 0.575 & 17.128 & 0.528\\

\bottomrule\end{longtable}\endgroup

\begingroup\small\setlength{\tabcolsep}{3pt}\begin{longtable}{llrrrrrr}\caption{Перебор селектора: duration.}\\\toprule

Process & Threshold & D & R & Agree & Cost ratio & Regret & Event\\\midrule\endfirsthead

\toprule Process & Threshold & D & R & Agree & Cost ratio & Regret & Event\\\midrule\endhead

P1 & 0 & 3 & 3 & 0.250 & 0.542 & 2.765 & 0.378\\

P1 & 0.0025 & 1 & 3 & 0.250 & 0.548 & 2.895 & 0.384\\

P1 & 0.005 & 0 & 3 & 0.250 & 0.549 & 2.920 & 0.384\\

P1 & 0.01 & 0 & 3 & 0.250 & 0.549 & 2.920 & 0.384\\

P1 & 0.02 & 0 & 3 & 0.250 & 0.549 & 2.920 & 0.384\\

P2 & 0 & 4 & 6 & 0.583 & 0.596 & 4.266 & 0.344\\

P2 & 0.0025 & 0 & 6 & 0.583 & 0.614 & 4.694 & 0.349\\

P2 & 0.005 & 0 & 6 & 0.583 & 0.614 & 4.694 & 0.349\\

P2 & 0.01 & 0 & 6 & 0.583 & 0.614 & 4.694 & 0.349\\

P2 & 0.02 & 0 & 6 & 0.583 & 0.614 & 4.694 & 0.349\\

P3 & 0 & 5 & 0 & 0.333 & 0.612 & 10.592 & 0.641\\

P3 & 0.0025 & 0 & 0 & 0.250 & 0.630 & 12.473 & 0.651\\

P3 & 0.005 & 0 & 0 & 0.250 & 0.630 & 12.473 & 0.651\\

P3 & 0.01 & 0 & 0 & 0.250 & 0.630 & 12.473 & 0.651\\

P3 & 0.02 & 0 & 0 & 0.250 & 0.630 & 12.473 & 0.651\\

P4 & 0 & 4 & 1 & 0.917 & 0.528 & 0.953 & 0.436\\

P4 & 0.0025 & 1 & 1 & 0.917 & 0.528 & 0.953 & 0.436\\

P4 & 0.005 & 0 & 1 & 0.917 & 0.528 & 0.953 & 0.436\\

P4 & 0.01 & 0 & 1 & 0.917 & 0.528 & 0.953 & 0.436\\

P4 & 0.02 & 0 & 1 & 0.917 & 0.528 & 0.953 & 0.436\\

P5 & 0 & 4 & 6 & 0.500 & 0.726 & 17.393 & 0.427\\

P5 & 0.0025 & 0 & 6 & 0.500 & 0.727 & 17.412 & 0.427\\

P5 & 0.005 & 0 & 6 & 0.500 & 0.727 & 17.412 & 0.427\\

P5 & 0.01 & 0 & 6 & 0.500 & 0.727 & 17.412 & 0.427\\

P5 & 0.02 & 0 & 6 & 0.500 & 0.727 & 17.412 & 0.427\\

P6 & 0 & 2 & 0 & 0.000 & 0.576 & 17.168 & 0.526\\

P6 & 0.0025 & 0 & 0 & 0.000 & 0.575 & 17.128 & 0.528\\

P6 & 0.005 & 0 & 0 & 0.000 & 0.575 & 17.128 & 0.528\\

P6 & 0.01 & 0 & 0 & 0.000 & 0.575 & 17.128 & 0.528\\

P6 & 0.02 & 0 & 0 & 0.000 & 0.575 & 17.128 & 0.528\\

\bottomrule\end{longtable}\endgroup

Исходные архивы траекторий содержат все независимые поступления и исходы кандидатных действий. Записи по обучающим наборам включают фактические пороги, их отношение к диапазону калибровки, активацию, время предупреждения, ложные и ненужные активации, интеграл WIP, длительности завершённых экземпляров и остаточную работу. Длительности завершённых экземпляров ограничены горизонтом оценки; проверки полной когорты, где доступны, являются отдельной диагностикой с продлённым горизонтом. Записи P4 сохраняют результаты обоих конкурирующих процессов.

\clearpage

\subsection{Диапазон наблюдений, активация и восстановление}

Таблица при нагрузочных возвратах показывает выбор retained-load. Числа порогов обозначают фактические значения WIP и число обучений. Диапазон обучения — разброс максимального наблюдавшегося WIP между наборами; $q/\max w$ считается внутри каждого набора. K/D обозначают известную/обнаруженную структуру. Диагностика не меняет правила активации.

\begin{center}\small\setlength{\tabcolsep}{3pt}\begin{tabular}{llp{.24\textwidth}rrr}\toprule

Process & N & Threshold:count & Train max & $q/\max w$ & Above\\\midrule

P1 & K & 6:11; 12:1 & 5--8 & 0.75--2.40 & 4/12\\

P1 & D & 6:12 & 5--8 & 0.75--1.20 & 4/12\\

P2 & K & 8:12 & 6--12 & 0.67--1.33 & 2/12\\

P2 & D & 8:12 & 6--12 & 0.67--1.33 & 2/12\\

P3 & K & 10:11; none:1 & 7--11 & 0.91--1.43 & 5/12\\

P3 & D & 10:11; none:1 & 7--11 & 0.91--1.43 & 5/12\\

P4 & K & 10:12 & 11--19 & 0.53--0.91 & 0/12\\

P4 & D & 10:12 & 11--19 & 0.53--0.91 & 0/12\\

P5 & K & 10:11; 20:1 & 12--18 & 0.56--1.18 & 1/12\\

P5 & D & 10:11; 20:1 & 12--18 & 0.56--1.18 & 1/12\\

P6 & K & 10:11; none:1 & 9--13 & 0.77--1.11 & 1/12\\

P6 & D & 10:12 & 9--13 & 0.77--1.11 & 1/12\\

\bottomrule\end{tabular}\end{center}

\begin{center}\small\setlength{\tabcolsep}{4pt}\begin{tabular}{llrrrrr}\toprule

Process & N & Activate & Nonsevere & Lead/$T$ & Event & Prevented\\\midrule

P1 & K & 0.717 & 0.371 & 5.661 & 0.365 & 0.326\\

P1 & D & 0.724 & 0.385 & 5.979 & 0.378 & 0.299\\

P2 & K & 0.736 & 0.542 & 5.666 & 0.325 & 0.223\\

P2 & D & 0.736 & 0.542 & 5.666 & 0.370 & 0.117\\

P3 & K & 0.628 & 0.052 & 4.649 & 0.651 & 0.023\\

P3 & D & 0.628 & 0.052 & 4.649 & 0.651 & 0.023\\

P4 & K & 0.828 & 0.533 & 8.362 & 0.490 & 0.223\\

P4 & D & 0.828 & 0.533 & 8.362 & 0.490 & 0.223\\

P5 & K & 0.833 & 0.640 & 9.790 & 0.354 & 0.330\\

P5 & D & 0.833 & 0.640 & 9.790 & 0.354 & 0.330\\

P6 & K & 0.667 & 0.203 & 6.346 & 0.401 & 0.382\\

P6 & D & 0.729 & 0.227 & 6.302 & 0.528 & 0.187\\

\bottomrule\end{tabular}\end{center}

Nonsevere — активация при отсутствии тяжёлого события в парной траектории без действия; это не обязательно вредное вмешательство. Время предупреждения измеряется относительно начала такого контрфактического события. Предотвращение условно по событию без действия. Восстановление ниже означает WIP не выше $2S$ в течение $10T$ после импульса с цензурированием на исходном горизонте $180T$. Сначала усредняются траектории каждого независимого обучения, затем сами обучения; учитываются все импульсные траектории, а не только кризисные.

\begin{center}\small\begin{tabular}{lrrrr}\toprule

Process & No action & Predictive K & Retained K & Generator\\\midrule

P1 & 0.792 & 1.000 & 1.000 & 1.000\\

P2 & 0.742 & 0.958 & 1.000 & 1.000\\

P3 & 0.169 & 0.667 & 0.560 & 0.727\\

P4 & 0.555 & 0.979 & 0.878 & 0.992\\

P5 & 0.815 & 0.917 & 0.987 & 1.000\\

P6 & 0.477 & 0.935 & 0.953 & 1.000\\

\bottomrule\end{tabular}\end{center}

\subsection{Последствия общих ресурсов и положение узкого места}

Ресурсная диагностика повторяет P4/P5 при всех четырёх механизмах и всех 12 обучениях для двух импульсных условий с высокой базовой интенсивностью и повтора 0. Включены все типы действий на минимальном фиксированном пороге и сохранённые выбранные политики. Все 880 префиксов до $180T$ точно совпадают с исходными тестовыми отсчётами. Затем поступления прекращаются, а действия сохраняются до $600T$ для завершения когорты. Интегралы очередей относятся только к исходному горизонту. Поэтому длительность полной когорты имеет другой горизонт, чем основной показатель.

\begin{longtable}{llrrrr}\caption{P4: интегралы WIP каждого процесса (тысячи) по полному тесту и времена завершения полной когорты по продлённой диагностике.}\\\toprule

Mechanism & Action at 10 & WIP X & WIP Y & Cycle X & Cycle Y\\\midrule\endfirsthead

\toprule Mechanism & Action at 10 & WIP X & WIP Y & Cycle X & Cycle Y\\\midrule\endhead

fixed & None & 7.78 & 8.61 & 21.53 & 23.86\\

fixed & Capacity & 4.68 & 5.41 & 11.94 & 13.80\\

fixed & Rework X & 7.44 & 8.45 & 20.61 & 23.36\\

fixed & Rework Y & 7.46 & 7.99 & 20.59 & 22.05\\

service & None & 43.69 & 45.54 & 184.34 & 191.92\\

service & Capacity & 22.81 & 24.08 & 84.37 & 87.91\\

service & Rework X & 41.85 & 44.62 & 174.80 & 186.61\\

service & Rework Y & 42.17 & 42.49 & 174.97 & 174.70\\

routing & None & 31.13 & 30.37 & 111.82 & 109.49\\

routing & Capacity & 15.01 & 14.88 & 46.73 & 46.50\\

routing & Rework X & 16.54 & 21.86 & 54.25 & 72.12\\

routing & Rework Y & 21.54 & 16.27 & 70.68 & 52.47\\

combined & None & 91.64 & 89.79 & 567.45 & 554.47\\

combined & Capacity & 72.57 & 70.68 & 373.43 & 365.29\\

combined & Rework X & 69.66 & 83.69 & 355.70 & 471.85\\

combined & Rework Y & 86.18 & 69.81 & 491.10 & 366.32\\

\bottomrule\end{longtable}

\begin{longtable}{llrrr}\caption{P5: средние очереди B и D на исходном горизонте; число парных траекторий с перемещением наибольшей очереди от B к D.}\\\toprule

Mechanism & Action at 10 & Queue B & Queue D & B to D\\\midrule\endfirsthead

\toprule Mechanism & Action at 10 & Queue B & Queue D & B to D\\\midrule\endhead

fixed & None & 13.12 & 0.11 & 0/24\\

fixed & Capacity B & 1.99 & 4.75 & 24/24\\

fixed & Capacity D & 13.12 & 0.01 & 0/24\\

fixed & Rework & 13.12 & 0.08 & 0/24\\

service & None & 140.94 & 0.01 & 0/24\\

service & Capacity B & 38.36 & 0.09 & 0/24\\

service & Capacity D & 140.94 & 0.00 & 0/24\\

service & Rework & 140.94 & 0.00 & 0/24\\

routing & None & 13.28 & 35.73 & 0/24\\

routing & Capacity B & 2.09 & 46.05 & 0/24\\

routing & Capacity D & 13.28 & 4.24 & 0/24\\

routing & Rework & 13.28 & 1.69 & 0/24\\

combined & None & 141.02 & 1.49 & 0/24\\

combined & Capacity B & 40.19 & 118.56 & 24/24\\

combined & Capacity D & 141.02 & 0.05 & 0/24\\

combined & Rework & 141.02 & 0.03 & 0/24\\

\bottomrule\end{longtable}

